\section*{6 Surface and Volume Integrals: From Flux to Global Theorems}

\subsection*{6.1 Surface Integrals}
\begin{conceptbox}
    The surface integral of a vector field over a surface $S$ is:
    \begin{equation*}
        \iint_S \mathbf{F} \cdot d\mathbf{S}
    \end{equation*}
\end{conceptbox}

This measures the \textbf{flux} of a vector field through a surface.

A surface is parametrized as:
\begin{equation*}
    \mathbf{r}(u,v)
\end{equation*}

Then:
\begin{equation*}
    d\mathbf{S} = (\mathbf{r}_u \times \mathbf{r}_v)\, dA
\end{equation*}

\begin{examplebox}
    Compute the flux of $\mathbf{F}=\langle x,y,z \rangle$ through the surface $z=4-x^2-y^2$, $z\ge0$.


    \textbf{Step 1: Understand the surface and determine its projection}

    The equation of the surface is
    \begin{equation*}
        z=4-x^2-y^2
    \end{equation*}

    Because the surface satisfies $z\ge 0$, we impose
    \begin{align*}
        4-x^2-y^2 & \ge 0 \\
        x^2+y^2   & \le 4
    \end{align*}

    So the projection of the surface onto the $xy$-plane is the disk
    \begin{equation*}
        x^2+y^2 \le 4
    \end{equation*}
    which is a circle of radius $2$ centered at the origin.

    This tells us the region of integration in the plane.

    \textbf{Step 2: Parametrize the surface}

    Since the surface is already given as $z$ in terms of $x$ and $y$, the natural parametrization is
    \begin{equation*}
        \mathbf{r}(x,y)=\langle x,\; y,\; 4-x^2-y^2\rangle
    \end{equation*}

    Here:
    \begin{itemize}
        \item the first component is the $x$-coordinate
        \item the second component is the $y$-coordinate
        \item the third component is the corresponding $z$-value on the surface
    \end{itemize}

    So every point $(x,y)$ in the disk $x^2+y^2\le 4$ gives a point on the surface.

    \textbf{Step 3: Compute the partial derivatives of the parametrization}

    To find the oriented surface element, we compute the tangent vectors:
    \begin{align*}
        \mathbf{r}_x
         & = \frac{\partial}{\partial x}\langle x,\; y,\; 4-x^2-y^2\rangle                                                                       \\
         & = \left\langle \frac{\partial x}{\partial x},\; \frac{\partial y}{\partial x},\; \frac{\partial}{\partial x}(4-x^2-y^2) \right\rangle \\
         & = \langle 1,\; 0,\; -2x \rangle
    \end{align*}

    Similarly,
    \begin{align*}
        \mathbf{r}_y
         & = \frac{\partial}{\partial y}\langle x,\; y,\; 4-x^2-y^2\rangle                                                                       \\
         & = \left\langle \frac{\partial x}{\partial y},\; \frac{\partial y}{\partial y},\; \frac{\partial}{\partial y}(4-x^2-y^2) \right\rangle \\
         & = \langle 0,\; 1,\; -2y \rangle
    \end{align*}

    These two vectors lie tangent to the surface.

    \textbf{Step 4: Compute the cross product $\mathbf{r}_x \times \mathbf{r}_y$}

    The cross product gives a normal vector to the surface:
    \begin{equation*}
        \mathbf{r}_x \times \mathbf{r}_y =
        \begin{vmatrix}
            \mathbf{i} & \mathbf{j} & \mathbf{k} \\
            1          & 0          & -2x        \\
            0          & 1          & -2y
        \end{vmatrix}
    \end{equation*}

    Expand the determinant:

    \begin{align*}
        \mathbf{r}_x \times \mathbf{r}_y
         & =
        \mathbf{i}
        \begin{vmatrix}
            0 & -2x \\
            1 & -2y
        \end{vmatrix}
        -
        \mathbf{j}
        \begin{vmatrix}
            1 & -2x \\
            0 & -2y
        \end{vmatrix}
        +
        \mathbf{k}
        \begin{vmatrix}
            1 & 0 \\
            0 & 1
        \end{vmatrix}                               \\
         & =
        \mathbf{i}\bigl(0(-2y)-(-2x)(1)\bigr)
        -
        \mathbf{j}\bigl(1(-2y)-(-2x)(0)\bigr)
        +
        \mathbf{k}(1\cdot 1-0\cdot 0)                \\
         & =
        \mathbf{i}(2x)-\mathbf{j}(-2y)+\mathbf{k}(1) \\
         & =
        \langle 2x,\; 2y,\; 1\rangle
    \end{align*}

    So the vector surface element is
    \begin{equation*}
        d\mathbf{S} = (\mathbf{r}_x \times \mathbf{r}_y)\, dA
        = \langle 2x,\; 2y,\; 1\rangle dA
    \end{equation*}

    This normal points upward because its $k$-component is positive.

    \textbf{Step 5: Evaluate the vector field on the surface}

    The field is
    \begin{equation*}
        \mathbf{F}(x,y,z)=\langle x,y,z\rangle
    \end{equation*}

    But on the surface, $z=4-x^2-y^2$. So substitute this into $\mathbf{F}$:
    \begin{align*}
        \mathbf{F}(\mathbf{r}(x,y))
         & = \left\langle x,\; y,\; 4-x^2-y^2 \right\rangle
    \end{align*}

    \textbf{Step 6: Compute the dot product $\mathbf{F}\cdot (\mathbf{r}_x \times \mathbf{r}_y)$}

    The flux integrand is:
    \begin{align*}
        \mathbf{F}\cdot (\mathbf{r}_x \times \mathbf{r}_y)
         & =
        \left\langle x,\; y,\; 4-x^2-y^2 \right\rangle
        \cdot
        \langle 2x,\; 2y,\; 1\rangle    \\
         & = x(2x)+y(2y)+(4-x^2-y^2)(1) \\
         & = 2x^2+2y^2+4-x^2-y^2        \\
         & = x^2+y^2+4
    \end{align*}

    Therefore the flux integral becomes
    \begin{equation*}
        \iint_S \mathbf{F}\cdot d\mathbf{S}
        =
        \iint_R (x^2+y^2+4)\, dA
    \end{equation*}
    where $R$ is the disk $x^2+y^2\le 4$.

    \textbf{Step 7: Convert the integral to polar coordinates}

    Because the region $R$ is a disk, polar coordinates are the natural choice:
    \begin{align*}
        x       & = r\cos\theta    \\
        y       & = r\sin\theta    \\
        x^2+y^2 & = r^2            \\
        dA      & = r\,dr\,d\theta
    \end{align*}

    The disk $x^2+y^2\le 4$ becomes
    \begin{align*}
        0 \le r \le 2, \qquad 0 \le \theta \le 2\pi
    \end{align*}

    Substitute into the integral:
    \begin{align*}
        \iint_R (x^2+y^2+4)\, dA
         & =
        \int_0^{2\pi}\int_0^2 (r^2+4)\,r\,dr\,d\theta
    \end{align*}

    \textbf{Step 8: Simplify the integrand}

    Distribute the $r$:
    \begin{align*}
        (r^2+4)r = r^3+4r
    \end{align*}

    So:
    \begin{align*}
        \int_0^{2\pi}\int_0^2 (r^3+4r)\,dr\,d\theta
    \end{align*}

    \textbf{Step 9: Integrate with respect to $r$}

    \begin{align*}
        \int_0^2 (r^3+4r)\,dr
         & =
        \left[\frac{r^4}{4}+2r^2\right]_0^2                \\
         & =
        \left(\frac{2^4}{4}+2(2^2)\right)-\left(0+0\right) \\
         & =
        \left(\frac{16}{4}+8\right)                        \\
         & =
        4+8                                                \\
         & =
        12
    \end{align*}

    \textbf{Step 10: Integrate with respect to $\theta$}

    \begin{align*}
        \int_0^{2\pi} 12\, d\theta
         & = 12[\theta]_0^{2\pi} \\
         & = 12(2\pi)            \\
         & = 24\pi
    \end{align*}

    \textbf{Final Answer:}
    \begin{equation*}
        \boxed{24\pi}
    \end{equation*}
\end{examplebox}
